% =============================================================================
% Kapitel 1: Einleitung und Problemstellung
% =============================================================================

\section{Einleitung und Problemstellung}
\label{sec:einleitung}

ROSI (\textit{Realtime Optical Surface Inspection}) ist ein von der Grenzebach BSH GmbH
entwickeltes Inline-Scan-System zur automatisierten Qualitätskontrolle von Baustoffen.
Im laufenden Produktionsprozess werden Materialoberflächen erfasst, Defekte per
KI-Inferenz erkannt und das Material anschließend regelbasiert in Qualitätsstufen
klassifiziert. Der Betrieb erfolgt On-Premise direkt in der Produktionsanlage des
Kunden im 24/7-Dauerbetrieb.

Technisch ist ROSI als Python-Multiprocessing-System aufgebaut. Mehrere langlebige
Prozesse verarbeiten die Daten in einer Shared-Memory-basierten Pipeline, ergänzt durch
Queues und flüchtigen Zwischenspeicher in Form eines tmpfs-RAM-Dateisystems. Große
Bilddaten wandern dabei nicht durch Queues, sondern verbleiben in vom Betriebssystem
verwalteten Speicherblöcken; die Queues transportieren ausschließlich Referenzen. Die
in dieser Fallstudie untersuchten Fehlerklassen setzen genau an dieser Prozess- und
Speicherarchitektur an.

ROSI befindet sich in aktiver Entwicklung; die erste Kundeninbetriebnahme ist für Mitte
bis Ende 2026 geplant. Mit dem bevorstehenden Produktiveinsatz steigen die Anforderungen
an Verfügbarkeit und Fehlertoleranz: Ausfälle können Produktionsstillstände
verursachen, Inspektionsergebnisse verhindern und im ungünstigsten Fall Qualitätsmängel
unentdeckt lassen.

\subsection{Problemstellung}
\label{sec:problemstellung}

Systeme im 24/7-Betrieb sind Fehlerklassen ausgesetzt, die im Entwicklungsbetrieb
unsichtbar bleiben. Ein System, das täglich neu gestartet wird und unter manueller
Beobachtung steht, durchläuft selten jene Zustände, in denen sich Ressourcen schleichend
erschöpfen, Verbindungen veralten oder Teilprozesse still ausfallen. Für ROSI lassen
sich die relevanten Risikokategorien wie folgt einteilen:

\begin{itemize}
    \item \textit{Speichererschöpfung:} Puffer, Logs oder Datenbanken wachsen
          schleichend, bis Schreiboperationen systemweit fehlschlagen
          (Abschnitt~\ref{sec:grundlagen-ressourcen}).
    \item \textit{Stille Prozessausfälle:} Ein Teilprozess stürzt ab, ohne dass das
          System dies erkennt (Fail-Silent-Verhalten, Abschnitt~\ref{sec:grundlagen-fehler}).
    \item \textit{Verbindungsabbrüche:} Datenbankverbindungen veralten nach
          Produktionspausen und führen beim nächsten Zugriff zu unkontrollierten
          Prozessabstürzen (Abschnitt~\ref{sec:grundlagen-parity}).
    \item \textit{Kaskadeneffekte:} Die Verlangsamung einer Komponente blockiert über
          gemeinsame Queues und Shared-Memory-Schnittstellen das gesamte System
          (Abschnitt~\ref{sec:grundlagen-fehler}).
    \item \textit{Blinde Flecken im Monitoring:} Vorhandene Überwachungskomponenten
          erfassen nicht alle relevanten Ressourcen, etwa flüchtige RAM-Dateisysteme
          oder interne Queue-Tiefen.
\end{itemize}

ROSI verfügt aktuell über kein externes Alerting. Probleme werden erst durch
ausbleibende Inspektionsergebnisse oder manuelle Beobachtung sichtbar -- die
Konsequenzen dieser Lücke werden in Kapitel~\ref{sec:fazit} quantitativ eingeordnet.

\subsection{Zielsetzung und Forschungsfragen}
\label{sec:zielsetzung}

Ziel dieser Fallstudie ist es, verfügbarkeitskritische Fehlerszenarien im industriellen
Dauerbetrieb auf Prozess- und Betriebssystemebene zu identifizieren, reproduzierbar zu
simulieren und anhand messbarer Kriterien systematisch zu bewerten. Im Fokus stehen
stille, schleichende und kaskadierende Fehler. Die untersuchten Failure Cases werden in
drei Themenblöcke gegliedert, denen jeweils eine Forschungsfrage (FF) zugeordnet ist:

\begin{itemize}
    \item \textit{FF1 -- Speicher- und Ressourcenverwaltung:} Wie verhält sich das
          System, wenn flüchtige Speicher\-subsysteme überlaufen oder Timing-Konflikte
          im Shared-Memory-Management auftreten?
    \item \textit{FF2 -- Systembeobachtbarkeit und Fehlertoleranz:} In welchem Zeitraum
          erkennt das System eigene Prozessausfälle, und welche Reaktion erfolgt?
    \item \textit{FF3 -- Datenpersistenz und Kaskadeneffekte:} Wie verhält sich das
          System bei unterbrochenen Datenbankverbindungen oder blockierenden Queues,
          und welche Folgeausfälle entstehen?
\end{itemize}

\subsection{Aufbau der Arbeit}
\label{sec:aufbau}

Kapitel~\ref{sec:architektur} beschreibt die für die Fehleranalyse relevante Prozess- und
Speicherarchitektur von ROSI. Kapitel~\ref{sec:grundlagen} legt den theoretischen Rahmen zu
Verfügbarkeit, Fehlerklassifikation, Prozesssupervision und Ressourcenmanagement.
Kapitel~\ref{sec:methodik} beschreibt die Testumgebung und die Methodik der Fehlerinjektion.
Kapitel~\ref{sec:cases} stellt die fünf Failure Cases mit ihren Messergebnissen dar.
Kapitel~\ref{sec:fazit} fasst die Befunde zusammen und leitet Handlungsempfehlungen ab.
